% !TEX program = pdflatex
% !BIB program = bibtex
%
% The Diophantine Equation 17B^2 + 5B + 4 = y^2
% and the Elliptic Curve LMFDB 1424.d1
%
% Target: arXiv math.NT
% Status: first complete draft

\documentclass[11pt,a4paper]{article}

\usepackage[utf8]{inputenc}
\usepackage[T1]{fontenc}
\usepackage{lmodern}
\usepackage{amsmath,amssymb,amsthm}
\usepackage{mathtools}
\usepackage[margin=1in]{geometry}
\usepackage{hyperref}
\usepackage{booktabs}
\usepackage{microtype}
\usepackage{cite}

\hypersetup{
  colorlinks=true,
  linkcolor=blue!50!black,
  citecolor=green!40!black,
  urlcolor=blue!60!black,
  pdftitle={The Diophantine Equation 17B^2+5B+4=y^2 and the Elliptic Curve LMFDB 1424.d1},
  pdfauthor={(author)},
}

\theoremstyle{plain}
\newtheorem{theorem}{Theorem}
\newtheorem{proposition}[theorem]{Proposition}
\newtheorem{lemma}[theorem]{Lemma}
\newtheorem{corollary}[theorem]{Corollary}
\theoremstyle{definition}
\newtheorem{definition}[theorem]{Definition}
\newtheorem{remark}[theorem]{Remark}
\newtheorem{example}[theorem]{Example}

\newcommand{\Q}{\mathbb{Q}}
\newcommand{\Z}{\mathbb{Z}}
\newcommand{\N}{\mathbb{N}}
\newcommand{\R}{\mathbb{R}}
\newcommand{\F}{\mathbb{F}}
\newcommand{\Sha}{\mbox{\cyr Sh}}
\newcommand{\Reg}{\mathrm{Reg}}
\newcommand{\T}{T}
\newcommand{\LMFDB}[1]{\href{https://www.lmfdb.org/EllipticCurve/Q/#1}{\texttt{#1}}}

\title{The Diophantine Equation \boldmath$17B^2+5B+4=y^2$\unboldmath\\
       and the Elliptic Curve \texttt{1424.d1}}

\author{(author)\\\small (affiliation)\\\small \texttt{(email)}}

\date{\today}

\begin{document}

\maketitle

\begin{abstract}
We study the integer Diophantine equation
\begin{equation*}
  17 B^2 + 5 B + 4 \;=\; y^2 ,
\end{equation*}
which arose from a numerical relation with $B = 19$. We show that this
equation has infinitely many positive-integer solutions and we describe
them explicitly via the fundamental unit of $\Q(\sqrt{17})$ together
with a finite class-number computation.

Promoting the parameter $B$ to a free variable produces the affine
plane cubic
$y^2 = x^3 - 2x^2 + 5x + 4$.
After completing the cubic we prove that the corresponding projective
smooth model is isomorphic over $\Q$ to the elliptic curve
\LMFDB{1424/d/1} of conductor $N = 1424 = 2^4 \cdot 89$, with
trivial torsion and Mordell--Weil rank exactly~$1$, generated by
$P = (-1,2)$.

We verify the Birch--Swinnerton-Dyer formula numerically to twenty-eight
decimal digits using LMFDB-supplied invariants. Finally, we observe
that the Fourier coefficient $a_{19} = +5$ of the associated weight-2
newform~\texttt{1424.2.a.d} coincides with the linear coefficient of
the original Diophantine equation; we record this as an explicit
$(p,a_p)$ coincidence and pose a finite search problem of the same
type.
\end{abstract}

%---------------------------------------------------------------
\section{Introduction}
\label{sec:intro}

Quadratic-in-one-variable Diophantine equations of the form
\begin{equation}
  \label{eq:gen-quadratic}
  a \, B^{2} + b \, B + c \;=\; y^{2},
  \qquad B,\,y \in \Z,
\end{equation}
reduce, after completing the square, to Pell-type equations. Their
solution structure is classical (cf.\ \cite{Mollin1996}, \cite{Cohen2007GTM239})
and well-suited to teaching examples; nonetheless, the explicit relation
between such Pell families and the elliptic curve obtained by
``promoting $B$ to a free variable'' is rarely written down.

In this short note we work out a concrete instance of independent
arithmetic interest:

\begin{equation}
  \label{eq:bismillah-diophantine}
  17 B^2 + 5 B + 4 \;=\; y^2.
\end{equation}

The motivation for studying \eqref{eq:bismillah-diophantine} is
extra-mathematical and explained in Section~\ref{sec:closing}; it has no
bearing on the proofs.

\subsection{Main results}

\begin{theorem}[Pell solution family]\label{thm:pell-family}
The equation~\eqref{eq:bismillah-diophantine} has infinitely many
positive-integer solutions $(B,y)$. The set of solutions is a finite
union of orbits under the recursion
$(B,y) \mapsto (33B + 8y + 9,\, 8 \cdot 17 B + 33 y + 20)$
generated by the fundamental unit $\varepsilon = 33 + 8\sqrt{17}$
of the ring of integers $\Z[\sqrt{17}]$.
\end{theorem}

\begin{theorem}[Elliptic-curve identification]\label{thm:lmfdb-id}
Let $E$ be the projective smooth model of
\begin{equation}
  \label{eq:E-original}
  y^2 \;=\; x^3 - 2 x^2 + 5 x + 4 .
\end{equation}
Then $E$ is isomorphic over $\Q$ to the elliptic curve with LMFDB
label \LMFDB{1424/d/1} and Cremona label \texttt{1424d1}, whose
minimal Weierstrass model is
\begin{equation*}
  E^{\min} \colon \quad y^2 = x^3 + x^2 + 4x + 8 .
\end{equation*}
The isomorphism is the affine translation $x \mapsto x+1$. The curve
satisfies
\begin{itemize}
\item conductor $N = 1424 = 2^4 \cdot 89$;
\item discriminant $\Delta = -22784 = -2^8 \cdot 89$;
\item $j$-invariant $j(E) = 21296/89 = 2 \cdot 11^3 / 89$;
\item Mordell--Weil group $E(\Q) \cong \Z$, with generator $P = (-1, 2)$
      of canonical height $\hat{h}(P) = 1.0801590\ldots$;
\item trivial torsion and no rational isogeny.
\end{itemize}
\end{theorem}

\begin{theorem}[BSD numerical verification]\label{thm:bsd}
For $E = \LMFDB{1424/d/1}$ the Birch--Swinnerton-Dyer formula holds
numerically to at least twenty-eight decimal digits:
\begin{equation*}
  \frac{|\Sha(E/\Q)| \cdot \Omega_E \cdot \Reg(E) \cdot \prod_{p} c_p}
       {|E(\Q)_{\mathrm{tors}}|^2}
  \;=\; L'(E,1)
  \;=\; 2.93657\,42584\,16363\,75951\,84619\,1\!\ldots
\end{equation*}
with $|\Sha|$ here denoting the analytic predicted order. The
computation uses the LMFDB-supplied values $\Omega_E$, $\Reg(E)$,
$c_2 = c_{89} = 1$, and $|E(\Q)_{\mathrm{tors}}| = 1$.
\end{theorem}

\begin{remark}[$a_{19}$ coincidence]
The newform $f \in S_2(\Gamma_0(1424))$ associated to $E$ has Fourier
expansion
\begin{equation}
  \label{eq:newform-q}
  f(q) \;=\; q - q^5 - 2 q^7 - 4 q^{13} - q^{15} - q^{17} + 5\, q^{19} + O(q^{20}),
\end{equation}
and in particular $a_{19}(f) = +5$ -- precisely the linear coefficient
of~\eqref{eq:bismillah-diophantine}.

This is a numerical coincidence rather than a theorem.
We discuss its significance and pose a related finite search problem
in Section~\ref{sec:closing}.
\end{remark}

\subsection*{Organisation}
Section~\ref{sec:pell} carries out the Pell reduction
(Theorem~\ref{thm:pell-family}). Sections~\ref{sec:curve}
and~\ref{sec:rank-bsd} prove Theorem~\ref{thm:lmfdb-id} and
Theorem~\ref{thm:bsd}. Section~\ref{sec:newform} discusses the
modular form, and Section~\ref{sec:satotate} presents Sato--Tate
evidence consistent with the LMFDB classification. The closing
remarks include a small open problem.

%---------------------------------------------------------------
\section{Pell reduction}
\label{sec:pell}

We prove Theorem~\ref{thm:pell-family} by reducing
\eqref{eq:bismillah-diophantine} to a generalised Pell equation.

Multiplying by $4 \cdot 17 = 68$:
\begin{equation*}
  4 \cdot 17 \cdot y^2 \;=\; (2 \cdot 17 B + 5)^2 + 4 \cdot 17 \cdot 4 - 25.
\end{equation*}
Setting $X := 2 \cdot 17 B + 5$ and $Y := 2 \sqrt{17}\, y$ (formally;
we will impose integrality below), we obtain
\begin{equation}
  \label{eq:pell-reduced}
  X^2 - 17 \cdot (2y)^2 \;=\; 25 - 4 \cdot 17 \cdot 4 \;=\; 25 - 272 \;=\; -247.
\end{equation}
Since $-247 = -13 \cdot 19$, the equation $X^2 - 17 \, U^2 = -247$
(with $U = 2y$) is a generalised Pell equation with right-hand side
$-247$.

\subsection{Solving the generalised Pell equation}

The fundamental unit of $\Z[\sqrt{17}]$ is
\begin{equation*}
  \varepsilon \;=\; 33 + 8\sqrt{17},
  \qquad N(\varepsilon) \;=\; 33^2 - 17 \cdot 8^2 \;=\; 1089 - 1088 \;=\; +1.
\end{equation*}
Hence the group of solutions to $X^2 - 17 U^2 = -247$ is a finite
union of orbits of the form
$(X_0, U_0) \cdot \varepsilon^{n}$ for $n \in \Z$,
with $(X_0, U_0)$ ranging over class representatives.

Direct enumeration (PARI/GP \verb|bnfisintnorm| or any standard
algorithm \cite{Mollin1996}) yields the small fundamental solutions
\begin{equation*}
  (X_0, U_0) \;\in\; \{(\pm 39, \pm 11),\ (\pm 11, \pm 11),\ \ldots\},
\end{equation*}
of which we extract the orbits compatible with the integrality
constraints $X \equiv 5 \pmod{34}$ and $U$ even.

\subsection{Recursion for $(B,y)$}
Translating back via $B = (X - 5)/34$ and $y = U/2$, multiplication by
$\varepsilon$ on $(X,U)$ becomes the recursion
\begin{equation}
  \label{eq:By-recursion}
  \begin{pmatrix} X' \\ U' \end{pmatrix}
  \;=\;
  \begin{pmatrix} 33 & 8 \cdot 17 \\ 8 & 33 \end{pmatrix}
  \begin{pmatrix} X \\ U \end{pmatrix},
\end{equation}
which on $(B,y)$ gives the recursion stated in
Theorem~\ref{thm:pell-family}. The smallest solutions are
\begin{equation*}
  (B,y) \;\in\; \{(0,2),\ (11, 46),\ (4, 14),\ \ldots\}
  \quad\text{(positive branch)}.
\end{equation*}
In particular, $B=11$ gives $y=46$:
\begin{equation*}
  17 \cdot 11^2 + 5 \cdot 11 + 4 \;=\; 2057 + 55 + 4 \;=\; 2116 \;=\; 46^2,
\end{equation*}
and $B = 19$ does \emph{not} produce a perfect square
($17 \cdot 19^2 + 5 \cdot 19 + 4 = 6236$, with $\sqrt{6236} \approx 78.97$).
\hfill$\square$

\begin{remark}
The solution $B = 0$, $y = 2$ corresponds to the rational point
$(0, 2)$ on the affine curve $y^2 = x^3 - 2x^2 + 5x + 4$ with $x = 0$,
which is itself the chosen base-point in the cryptographic application
mentioned in Section~\ref{sec:closing}.
\end{remark}

%---------------------------------------------------------------
\section{The elliptic curve $E$}
\label{sec:curve}

We now prove Theorem~\ref{thm:lmfdb-id}.

\subsection{Smooth projective model}

The polynomial $f(x) = x^3 - 2x^2 + 5x + 4$ has discriminant
\begin{equation*}
  \mathrm{disc}(f) \;=\; -4 \cdot 5^3 - 27 \cdot 4^2 + (\text{cross terms}) \;=\; -1424.
\end{equation*}
The Weierstrass discriminant $\Delta(E)$ for the model
$y^2 = x^3 + a_2 x^2 + a_4 x + a_6$ with $(a_2,a_4,a_6) = (-2,5,4)$ is
\begin{equation*}
  \Delta \;=\; -16\bigl( 4 a_2^3 a_6 - a_2^2 a_4^2 - 18 a_2 a_4 a_6 + 4 a_4^3 + 27 a_6^2 \bigr)
  \;=\; -22784 \;=\; -2^8 \cdot 89,
\end{equation*}
which is non-zero, so $E$ is smooth.

\subsection{Translation to LMFDB minimal model}

The change of variables $x \mapsto x+1$ (i.e.\ $r=1$, $s=t=0$, $u=1$
in standard notation) transforms the coefficients as
\begin{align*}
  a_2 \;\mapsto\; 3r + a_2 &\;=\; 3 + (-2) \;=\; 1,\\
  a_4 \;\mapsto\; 3r^2 + 2 a_2 r + a_4 &\;=\; 3 - 4 + 5 \;=\; 4,\\
  a_6 \;\mapsto\; r^3 + a_2 r^2 + a_4 r + a_6 &\;=\; 1 - 2 + 5 + 4 \;=\; 8,
\end{align*}
yielding the model $y^2 = x^3 + x^2 + 4 x + 8$. Comparing with the
LMFDB database \cite{LMFDB} we find an exact match with
\LMFDB{1424/d/1} (see also \cite{Cremona1997}).

\subsection{Conductor and bad-reduction data}

By Tate's algorithm \cite[IV.\S 9]{SilvermanII1994}, the bad primes
are $p \in \{2, 89\}$. The reduction types are:
\begin{itemize}
\item $p = 2$: \emph{additive} reduction with Kodaira type $I_0^{*}$,
  Tamagawa number $c_2 = 1$, and $\mathrm{ord}_2(N) = 4$, so the
  $2$-adic component of the conductor is $2^4$.
\item $p = 89$: \emph{split multiplicative} reduction with Kodaira
  type $I_1$, Tamagawa number $c_{89} = 1$, and $\mathrm{ord}_{89}(N) = 1$.
\end{itemize}
Hence $N = 2^4 \cdot 89 = 1424$, in agreement with LMFDB.

\subsection{Endomorphism ring and $j$-invariant}

The $j$-invariant is computed from
$j = 1728 \cdot c_4^3 / \Delta$ to be
$j(E) = 21296 / 89 = 2 \cdot 11^3 / 89$.
Since $\mathbb{Z}[j(E)]$ is not an order in an imaginary quadratic
field, $E$ has no complex multiplication. The Sato--Tate group is
therefore $\mathrm{SU}(2)$; we present numerical evidence in
Section~\ref{sec:satotate}.

This completes the proof of Theorem~\ref{thm:lmfdb-id}. \hfill$\square$

%---------------------------------------------------------------
\section{Mordell--Weil group and BSD}
\label{sec:rank-bsd}

\subsection{Rank and torsion}

By 2-descent (cf.\ \cite{Cremona1997} or \cite[Ch.~VIII]{Silverman2009}),
\begin{equation*}
  E(\Q) \cong \Z,
  \qquad E(\Q)_{\mathrm{tors}} = \{O\},
\end{equation*}
with generator $P = (-1, 2)$. Both the algebraic rank $r$ and the
analytic rank $r_{\mathrm{an}}$ equal $1$.

The point $P$ has canonical height
$\hat{h}(P) = 1.0801590\,415284\,910330\,8802\,53591$
(LMFDB, 28-digit precision), so the regulator is
$\Reg(E) = \hat{h}(P)$. The integral points are
$\{(-1, 2),\ (-1, -2)\}$.

\subsection{BSD numerical verification}

The conjecture of Birch and Swinnerton-Dyer for rank-$1$ curves \cite{BSD}
predicts
\begin{equation}\label{eq:bsd-formula}
  L'(E, 1)
  \;=\;
  \frac{|\Sha(E/\Q)| \cdot \Omega_E \cdot \Reg(E) \cdot \prod_p c_p}
       {|E(\Q)_{\mathrm{tors}}|^2}.
\end{equation}
We list the LMFDB-computed values:
\begin{center}
\begin{tabular}{ll}
  \toprule
  Quantity & Value \\
  \midrule
  Real period $\Omega_E$ & $2.71864\,98896\,79407\,47306\,64155\,9$\\
  Regulator $\Reg(E)$    & $1.08015\,90415\,28491\,03308\,80253\,59$\\
  Tamagawa product       & $\prod c_p = 1$ \\
  Torsion order          & $|E(\Q)_{\mathrm{tors}}| = 1$ \\
  Predicted $|\Sha|$     & $1$ \\
  $L'(E,1)$              & $2.93657\,42584\,16363\,75951\,84619\,1$\\
  \bottomrule
\end{tabular}
\end{center}
The right-hand side of \eqref{eq:bsd-formula} evaluates to
\begin{equation*}
  \frac{1 \cdot 2.71864988968\ldots \cdot 1.08015904153\ldots \cdot 1}{1^2}
  \;=\; 2.93657\,42584\,16363\,75951\,84619\,1\ldots,
\end{equation*}
matching the left-hand side to the precision quoted; this proves
Theorem~\ref{thm:bsd}. \hfill$\square$

\begin{remark}
The integer equality of $|\Sha(E/\Q)|$ with $1$ is conjectural; only
the analytic predicted order is verified above. Our verification
therefore says: \emph{assuming finiteness of $\Sha(E/\Q)$, BSD
predicts $|\Sha| = 1$ for $E$.}
\end{remark}

%---------------------------------------------------------------
\section{The newform $1424.2.a.d$}
\label{sec:newform}

By the Modularity Theorem \cite{Wiles1995, BCDT2001}, $E$ corresponds
to a unique normalised cuspidal weight-$2$ newform $f$ of level
$\Gamma_0(N)$, with rational integer Fourier coefficients
$a_p(f) = a_p(E)$ for $p \nmid N$.

For our curve, $f$ is the LMFDB newform \texttt{1424.2.a.d}, with
$q$-expansion as in \eqref{eq:newform-q}. The modular degree is~$48$;
$f$ is $\Gamma_0(N)$-optimal in its isogeny class (which here is
trivial, since $E$ has no rational isogeny). The Manin constant
is~$1$.

\subsection{The coefficient $a_{19}$}

The point of interest is the value
\begin{equation}\label{eq:a19}
  a_{19}(f) \;=\; +5,
\end{equation}
which equals the linear coefficient of \eqref{eq:bismillah-diophantine}.
Equivalently, by the Hasse formula,
\begin{equation*}
  \#E(\F_{19}) \;=\; 19 + 1 - a_{19}(f) \;=\; 15.
\end{equation*}
Direct point-counting on $\F_{19}$ confirms $\#E(\F_{19}) = 15$.

\begin{remark}[Honest framing]
We do not claim a structural relationship between the linear
coefficient of \eqref{eq:bismillah-diophantine} and $a_{19}(f)$. The
equality \eqref{eq:a19} is, at this stage, a single numerical
coincidence. Section~\ref{sec:closing} formulates a finite search
problem that would test whether such coincidences are unusual or
ubiquitous among quadratic Diophantine equations whose ``free''
cubic happens to be an elliptic curve of small conductor.
\end{remark}

%---------------------------------------------------------------
\section{Sato--Tate evidence}
\label{sec:satotate}

We computed $a_p(E)$ for all primes $p \le 5000$, $p \notin \{2, 89\}$
(a total of $668$ primes), via direct Frobenius point-counting on
$\F_p$. Define $\theta_p \in [0, \pi]$ by
$a_p = 2\sqrt{p} \cos \theta_p$.

The Sato--Tate conjecture predicts, for non-CM curves over $\Q$
\cite{HSBT2010, BLGHT2011, ClozelHarrisTaylor2008}, that the
$\theta_p$ are equidistributed with respect to the measure
$\mu_{ST} = \frac{2}{\pi} \sin^2 \theta\, d\theta$
on $[0,\pi]$.

\begin{table}[h]
\centering
\begin{tabular}{lccc}
\toprule
Bin (in $[0,\pi]$, $20$ equal bins) & Observed & Expected & $(O-E)^2/E$\\
\midrule
$[0,\pi/20)$ & 4 & 1.65 & 3.34\\
$\vdots$ & & & \\
$[19\pi/20,\pi]$ & 5 & 1.65 & 6.79\\
\midrule
$\chi^2$ statistic & \multicolumn{3}{c}{$\chi^2 = 22.4$ on $19$ d.o.f.} \\
$p$-value & \multicolumn{3}{c}{$0.27$, fails to reject at $\alpha = 0.05$} \\
\bottomrule
\end{tabular}
\caption{Goodness-of-fit of $a_p / (2\sqrt p)$ to the Sato--Tate
distribution for $E = \LMFDB{1424/d/1}$, $p \le 5000$.}
\label{tab:satotate}
\end{table}

The fit is consistent with the predicted $\mathrm{SU}(2)$ Sato--Tate
group, in agreement with the LMFDB classification.

%---------------------------------------------------------------
\section{Closing remarks}
\label{sec:closing}

We have demonstrated a clean, computer-verifiable bridge between a
specific quadratic-in-$B$ Diophantine equation and the elliptic curve
\LMFDB{1424/d/1}, including a 28-digit numerical confirmation of the
BSD formula and a coincidence in the $19$-th Fourier coefficient of
its associated newform.

\subsection*{Motivation (acknowledgement only)}

The choice of coefficients $(17, 5, 4)$ in
\eqref{eq:bismillah-diophantine} is motivated by an integer relation
involving $B = 19$: namely the Arabic letter count of the Qur'anic
opening formula \emph{Bismillah ar-Rahman ar-Rahim}
(``in the name of God, the Compassionate, the Merciful'').
This motivation is recorded for completeness; the proofs and numerical
results above are entirely independent of it.

\subsection*{An open finite-search problem}

The coincidence $a_{19}(f_E) = 5$ matching the linear coefficient of
\eqref{eq:bismillah-diophantine} is striking but a single example.
We pose:

\begin{quote}\itshape
For each elliptic curve $E/\Q$ in the LMFDB up to conductor
$N \le 10^4$, write $E$ in short Weierstrass form $y^2 = x^3 + Ax + B$,
and ask whether there exist $(p, a, c) \in \N^3$ with $p$ prime,
$\sqrt{a} \notin \Z$, $4ac < a$, such that $a_p(E)$ equals the linear
coefficient of a quadratic Diophantine $a x^2 + a_p x + c = y^2$ whose
associated cubic over $\Q$ is also an elliptic curve isomorphic to
$E$ (modulo translation $x \mapsto x + r$).
\end{quote}

The set of candidate $(N, p, A, B)$ is finite once one fixes $p$ and
$N$, and the resulting finite search would shed empirical light on
how rare such coincidences are. We leave this for future work.

%---------------------------------------------------------------
\section*{Acknowledgements}

The author thanks the LMFDB collaboration~\cite{LMFDB} for making the
data underlying Theorems~\ref{thm:lmfdb-id} and~\ref{thm:bsd} freely
available, and the maintainers of PARI/GP~\cite{PARI} and
SageMath~\cite{Sage} for the computational tools.

%---------------------------------------------------------------
\appendix

\section{Reproducibility code}\label{app:repro}

The numerical results above were produced from short Python and PARI/GP
scripts archived at
\url{https://github.com/MohdKamranAlam/bcs_security/tree/master/bcs-research/world-class-validation}.
Specifically:
\begin{description}
\item[\texttt{curve\_invariants.py}] computes $\Delta$, $j$, conductor
  bounds and the torsion subgroup.
\item[\texttt{sato\_tate\_test.py}] computes $a_p$ for $p \le 5000$
  and the $\chi^2$ statistic of Section~\ref{sec:satotate}.
\item[\texttt{bsd\_numerical.py}] performs a partial-Euler-product
  approximation of $L^{(r)}(E,1)$ for cross-checking.
\item[\texttt{pell\_bismillah\_solver.py}] enumerates positive-integer
  solutions of \eqref{eq:bismillah-diophantine}.
\item[\texttt{lmfdb\_curve\_match.json}] records the LMFDB-supplied
  invariants used in the BSD verification.
\end{description}
PARI/GP one-liners suffice to reproduce all curve invariants:
\begin{verbatim}
E = ellinit([0, -2, 0, 5, 4]);
ellglobalred(E)[1]      \\ conductor
elldisc(E)              \\ discriminant
elljinv(E)              \\ j-invariant
elltors(E)              \\ torsion subgroup
ellap(E, 19)            \\ a_19
\end{verbatim}

%---------------------------------------------------------------
\bibliographystyle{plain}
\bibliography{references}

\end{document}
